\chapter{Conclusiones y trabajo futuro}
\label{ch:ch8_concl}

% [BORRADOR, 08-11-2026] Escrito antes de la puesta en marcha en la FPGA.
% Toda frase que depende de la placa va marcada [PENDIENTE] y hay que
% reescribirla con las cifras medidas (docs/puesta_en_marcha_fpga.md).

\section{Conclusiones}

El trabajo se propuso diseñar un motor de generación de claves para un
circuito integrado cuya entropía viene de osciladores en anillo simulados
a nivel de transistor, convertir su jitter en una cota mediante un modelo
estocástico, integrar la fuente en una \ac{FPGA} con un ESP32 y comparar
las secuencias aleatorias de ambos. Las conclusiones siguen los objetivos
de la Sección~\ref{sec:objectives}; cada una dice qué se ha demostrado y
con qué evidencia, \meas{}, \simu{} o \est{}.

\paragraph{La fuente de entropía se puede dimensionar desde el transistor
(O1--O3).} El anillo de cinco etapas en el proceso abierto IHP SG13G2
oscila a 2,73~GHz con un jitter de periodo de $176{,}7 \pm 2{,}0$~fs
\simu{} en la esquina típica a 27~\textdegree C, sacado del ruido de los
propios modelos de los dispositivos y no inyectado. El simulador no deja
fijar la realización de su fuente de ruido, así que la incertidumbre sale
de repetir la simulación, y la resolución del banco se declaró antes de
leer ninguna cifra. El cálculo exacto con la gaussiana envuelta, y no la
serie truncada que suele citarse, convierte el jitter en min-entropía por
bit, y el objetivo $H_\infty \ge 0{,}98$ fija $Q \ge 0{,}2286$. Entre
esquinas, alimentación ($\pm 10$\,\%) y desapareamiento, el peor caso es
la esquina lenta a $-40$~\textdegree C y 1,08~V; un divisor fijo de
$2^{23}$ cubre todos los chips simulados con $H_\infty \ge 0{,}9996$
\simu{}, a 0,20--0,41~kbit/s. El divisor sale, por tanto, de la esquina
fría, lenta y de alimentación baja, no de la típica, que por sí sola
habría sugerido $2^{20}$.

\paragraph{Una batería de estimadores estadísticos no sustituye al
modelo.} Bits sacados del modelo con el jitter de IHP, y bits uniformes
ideales, pasaron por la herramienta de SP~800-90B por la misma vía que las
capturas del ESP32. A $10^6$ bits la herramienta se satura en
$0{,}839 \pm 0{,}021$ \simu{}: todo punto del modelo con
$H_\infty \ge 0{,}97$ da lo mismo, y con poca entropía la herramienta
sobreestima. Con secuencias más largas el techo sube para todas las
fuentes, el ESP32 incluido, y la herramienta sigue sin separarlas
(Sección~\ref{sec:sim_90b}). Los tests estadísticos detectan defectos; no
certifican entropía. En un oscilador en anillo la cota tiene que venir de
la física, y ese es el argumento a favor del modelo estocástico.

\paragraph{El motor digital es correcto frente a referencias
independientes (O4).} El motor ---muestreador \ac{ERO}, tests de salud
con test de arranque, CTR\_DRBG con función de derivación, AES-CMAC,
zeroización y un mapa de registros por I2C--- está en VHDL-2008 portable.
Sus bancos de pruebas reproducen los vectores oficiales de FIPS-197,
SP~800-38B, RFC~4493 y \ac{CAVP}, y las reglas de seguridad del
secuenciador se comprueban como propiedades \ac{PSL} en toda la
simulación, no solo en los pasos que las provocan. Dos revisiones
independientes encontraron defectos que los bancos no habían cazado
---salida habilitada con una alarma de salud levantada, la semilla
llegando al buffer de captura, valores intermedios que sobrevivían a la
zeroización---; todos se corrigieron y cada uno tiene ahora su prueba. La
lección es de método: un diseño verificado por su autor contra su propio
banco no está verificado.

\paragraph{El motor cabe en un circuito integrado pequeño (O7).}
Sintetizado sobre las celdas estándar de IHP SG13G2, el motor ocupa
92,5~kGE con un único núcleo AES de 32 bits compartido \simu{}, de los que
unos 46~kGE son el motor digital sin el buffer de captura. Las dos
optimizaciones aplicadas ---un núcleo AES compartido por el generador y el
autenticador, y después una ruta de datos de 32 bits--- ahorraron
23~kGE; no se llega a los 20--25~kGE publicados en el estado del arte, y
está identificado el camino hasta unos 42~kGE.

\paragraph{El generador del ESP32 no es una referencia en la que confiar
a ciegas (O8, lado del ESP32).} Treinta y ocho capturas de tres chips en
cuatro condiciones de radio y \ac{ADC} no muestran efecto medible de la
condición ni del chip bajo SP~800-90B \meas{}. La prueba de reinicios
enseña lo que la estadística esconde: sin el ruido que mete el cargador de
arranque, el generador produce la misma secuencia tras cada reinicio, y
dos chips distintos producen la misma \meas{}. Su salida es tan buena como
su siembra, que depende de que la radio, el \ac{ADC} o el cargador de
arranque estén activos.

\paragraph{Integración en la FPGA (O5, O6).} [PENDIENTE: puesta en marcha
en la placa, frecuencias medidas de los anillos, $Q_{\mathrm{raw}}$ y
exponente de acumulación medidos, divisor para la \ac{FPGA}, campaña de
SP~800-90B con los anillos reales, reto-respuesta con el ESP32 en el bus.]
Antes de la placa, el firmware corrió contra un modelo del motor en el PC
con el mbedtls real, en un chip emulado y en un ESP32-C3 real; el
envoltorio de la placa, con su frecuencímetro de los anillos, se simuló
con los anillos del motor.

\paragraph{Comparación de las dos fuentes (O8).} [PENDIENTE: el motor
medido frente al ESP32 medido con la misma batería. Lo que anticipa la
simulación: la batería no los separará donde los dos están cerca de la
entropía plena; los puntos informativos son los divisores bajos, donde el
modelo predice $H_\infty < 0{,}8$ y la batería puede contradecirlo.]

\section{Aportaciones}

\begin{itemize}
\item Una cadena del ruido del transistor a la cota de min-entropía de un
oscilador en anillo sobre un \ac{PDK} abierto, con el divisor fijado por
el peor caso de esquinas, alimentación y desapareamiento, y con cada
eslabón validado con datos de respuesta conocida.
\item La evidencia, con fuentes simuladas de entropía conocida, de que los
estimadores de SP~800-90B se saturan por debajo de 0,9 a las longitudes
de secuencia habituales y no distinguen en ese rango una fuente sana de
una más débil.
\item Una caracterización del generador del ESP32 bajo SP~800-90B con tres
chips y cuatro condiciones, y una prueba de reinicios que muestra que, sin
ruido externo en el arranque, chips distintos producen la misma
secuencia. No se ha encontrado este resultado en la literatura
consultada.
\item Un motor de generación de claves en RTL portable con sus reglas de
seguridad escritas como propiedades comprobadas, sintetizado sobre celdas
estándar abiertas, con su firmware y un flujo reproducible publicados.
\end{itemize}

\section{Trabajo futuro}

\begin{itemize}
\item \emph{Silicio.} Fabricar el anillo y el muestreador en un shuttle
académico y medir el jitter que este trabajo ha simulado; la \ac{FPGA}
mide otro anillo.
\item \emph{Modelos de fundición.} Repetir el estudio de esquinas y
desapareamiento con los modelos de ruido de un \ac{PDK} comercial.
\item \emph{Captura continua.} El buffer de captura entrega bloques
independientes de 4096 bits; un modo continuo daría las muestras
consecutivas que pide SP~800-90B.
\item \emph{Test de jitter en línea.} Llevar el estimador de pares al RTL
como test en línea que detecte enfriamiento y enganche, como proponen sus
autores.
\item \emph{Área.} Aplicar la tercera optimización para acercarse a
42~kGE, y medir la potencia.
\item \emph{El producto.} Guardar la clave del equipo cifrada o hacer del
motor un elemento seguro que firme, y dar hora real a las lecturas
firmadas.
\item \emph{Verificación formal} de las propiedades del secuenciador,
cuando la cadena de herramientas las lea.
\end{itemize}
